\documentclass[10pt]{article}

\usepackage[margin=0.5in,top=0.35in,bottom=0.4in]{geometry}
\usepackage[T1]{fontenc}
\usepackage{lmodern}
\usepackage{microtype}
\usepackage{amssymb}
\usepackage{xcolor}
\usepackage{titlesec}
\usepackage{enumitem}
\usepackage[hidelinks]{hyperref}

\definecolor{accent}{HTML}{1F4E79}
\definecolor{contactgray}{HTML}{444444}

\pagestyle{empty}
\setlength{\parindent}{0pt}
\linespread{0.93}

% Section headers: bold colored title with a rule underneath
\titleformat{\section}
  {\large\bfseries\color{accent}}{}{0pt}{}
  [\vspace{1pt}\textcolor{accent}{\rule{\textwidth}{0.6pt}}]
\titlespacing*{\section}{0pt}{4.5pt}{2pt}

\setlist[itemize]{leftmargin=1.25em,itemsep=0.2pt,topsep=1pt,parsep=0pt}

\newcommand{\disc}[1]{{\fontsize{10.5}{12.6}\selectfont\textbf{#1}}}

\begin{document}

% ---------- Header ----------
\begin{center}
  {\Huge\bfseries Guang-Xing Li, Ph.D.}\\[1.5pt]
  {\large\color{accent}\bfseries Computational Physicist \,|\, Scientific Machine Learning \,|\, Multiscale Systems}\\[1.5pt]
  {\small\color{contactgray}
   \href{mailto:ligx.ngc7293@gmail.com}{ligx.ngc7293@gmail.com} \quad$|$\quad
   \href{https://github.com/gxli}{github.com/gxli} \quad$|$\quad
   \href{https://gxli.github.io}{gxli.github.io}}
\end{center}
\vspace{1.5pt}

% ---------- Summary ----------
\section*{Summary}
Computational physicist developing methods to extract structure, dynamics, and physical understanding
from complex scientific data. For 15+ years, his work has progressed from physics-based structural
diagnostics and multiscale representations to ScaleAware-JEPA, advancing self-supervised
representation learning toward \textbf{ML-assisted knowledge discovery directly from scientific
fields}.

{\small\textbf{Research scale:} \textbf{111 papers} $\cdot$ \textbf{1,800+ citations}
$\cdot$ \textbf{h-index 24} $\cdot$ \textbf{\texteuro280k+ PI-led funding} $\cdot$
international research leadership}

% ---------- Technical Skills ----------
% ---------- Selected ML Research ----------
\section*{Research Highlights --- Computational Methods for Scientific Discovery}

\href{https://arxiv.org/abs/2606.29723}{\textcolor{accent}{\textbf{ScaleAware-JEPA}}} ---
\textit{Representation learning for scientific knowledge discovery}

\hfill{\small\color{contactgray} 2026 $\cdot$ Sole author}\\[2pt]
{\small\color{contactgray} \href{https://arxiv.org/abs/2606.29723}{arXiv:2606.29723} $\cdot$ \href{https://github.com/gxli/SA-JEPA}{github.com/gxli/SA-JEPA}}\\[2pt]
A framework for using learned representations \textbf{not merely to predict predefined targets, but
to expose scientifically meaningful structure before the objects of interest are specified}.
\begin{itemize}
  \item Moves scientific ML \begingroup\setlength{\fboxsep}{1.2pt}\colorbox{accent!12}{\textbf{beyond predefined prediction toward knowledge
        discovery}}\endgroup: learns dense latent coordinates in which coherent structures can emerge
        \textbf{before objects, labels, or segmentation rules are specified}.
  \item Introduces \textbf{scale-aware latent prediction} for continuous physical fields, tying the
        JEPA prediction geometry to the intrinsic multiscale hierarchy of the data rather
        than arbitrary fixed image patches.
  \item Demonstrates a domain-general discovery framework across MHD turbulence,
        interstellar molecular gas, and urban structure; in astrophysical data, latent morphology
        recovers \textbf{physically distinct velocity--size relations}.
  \item Built the end-to-end research system, spanning multiscale
        representation, self-supervised GPU training, dense latent-field inference, and
        representation diagnostics.
\end{itemize}
\vspace{4pt}

\href{https://arxiv.org/abs/2506.05759}{\textbf{Adjacent Correlation Analysis (ACA)}} ---
\textit{Hidden regularities in complex spatial data}

\hfill{\small\color{contactgray} 2025 $\cdot$ Sole author}\\[2pt]
{\small\color{contactgray} \href{https://arxiv.org/abs/2506.05759}{arXiv:2506.05759} $\cdot$ \href{https://github.com/gxli/Adjacent-Correlation-Analysis}{github.com/gxli/Adjacent-Correlation-Analysis}}\\[2pt]
\begin{itemize}
  \item Introduced a spatially aware correlation framework that recovers \textbf{local relationships
        obscured by global statistics}, transforming complex spatial patterns into interpretable
        vector fields in phase space.
  \item Connects data-driven regularity discovery with dynamical-systems structure: the recovered
        correlation vectors can represent the \textbf{vector field on an attracting manifold} and
        expose \textbf{distinct physical regimes}.
\end{itemize}
\vspace{4pt}

\textbf{Constrained Diffusion Decomposition (CDD)} --- \textit{Multiscale representation of
continuous fields}

\hfill{\small\color{contactgray} 2022 $\cdot$ Sole author}\\[2pt]
{\small\href{https://arxiv.org/abs/2201.05484}{arXiv:2201.05484} $\cdot$ ApJS 259, 59 $\cdot$
\href{https://github.com/gxli/Constrained-Diffusion-Decomposition}{github.com/gxli/Constrained-Diffusion-Decomposition}}
\begin{itemize}
  \item Developed a nonlinear-diffusion decomposition that assigns \textbf{pixel-registered physical
        scales} while avoiding the ringing and negative artifacts introduced by conventional
        band-limited filtering.
  \item Provides the \textbf{explicit multiscale coordinate} later used in learned physical
        representations.
\end{itemize}
\vspace{4pt}

\textbf{G-virial} --- \textit{Physics-based identification of gravitational structure}

\hfill{\small\color{contactgray} 2014--2018 $\cdot$ DFG-funded}\\[2pt]
\begin{itemize}
  \item Developed an analytical--numerical framework for identifying \textbf{gravitationally
        significant structure} in turbulent media by connecting virial theory with simulated
        observables.
\end{itemize}
\vspace{4pt}

\section*{Research Leadership}

\textbf{Principal Investigator \& Associate Professor} \hfill \textbf{2018--Present}\\
\textit{Yunnan University (SWIFAR), China}\\
\begin{itemize}
  \item Lead a computational research program spanning multiscale physics, scientific ML, and
        analysis of observational and simulation data.
  \item Supervise PhD/MSc researchers in computational methods, scientific data analysis, and
        research software development.
  \item Advise collaborators on complex scientific datasets, translating physical questions into
        computational formulations, diagnostics, data-analysis pipelines, and ML approaches.
  \item \textbf{Led a \texteuro60k Max Planck Partner Group} supporting international collaborative
        research.
\end{itemize}

% ---------- Experience ----------
\newpage
\section*{Earlier Research Experience}

\textbf{DFG Research Fellow} \hfill \textbf{2014--2018}\\
\textit{LMU Munich (Universit\"ats-Sternwarte), Germany}\\
\begin{itemize}
  \item Awarded a \texteuro220,000 DFG Research Fellowship to develop computational methods for
        extracting physical structure and dynamical understanding from observational and
        simulation data.
  \item Developed G-virial, combining analytical physics and numerical simulations to identify
        gravitationally significant structure in turbulent molecular-gas data.
\end{itemize}

\textbf{Ph.D. Researcher} \hfill \textbf{2010--2014}\\
\textit{Max Planck Institute for Radio Astronomy, Germany}\\
\begin{itemize}
  \item Developed robust, scalable algorithms and computational models for multiscale structures in
        high-dimensional, heterogeneous environments; graduated \textit{magna cum laude}.
\end{itemize}

% ---------- Selected Scientific Publications ----------
\section*{Selected Scientific Discoveries}

\disc{A kiloparsec-scale, long-lived superbubble across the Perseus arm} --- \textit{Nature Communications} (2025)\\
Identified a large, long-lived, slowly expanding Galactic superbubble; \textbf{co-corresponding
author.}\\
B. Chen, \textbf{G.-X. Li}, H. Yuan, \textit{et al.}
\vspace{22pt}

\disc{Cosmic-ray shielding by dense molecular gas} --- \textit{Nature Astronomy} (2023)\\
Demonstrated effective suppression of $\lesssim$10\,GeV cosmic rays inside dense molecular clumps.\\
R.-Z. Yang, \textbf{G.-X. Li}, E. de O\~na Wilhelmi, \textit{et al.}
\vspace{22pt}

\disc{Flyby-induced spirals in a massive protostellar disk} --- \textit{Nature Astronomy} (2022)\\
Interpreted spiral structure in a massive Keplerian disk as the dynamical signature of a stellar
flyby; \textbf{co-corresponding author; led interpretation and analytical modeling.}\\
X. Lu, \textbf{G.-X. Li}, Q. Zhang \& Y. Lin.
\vspace{22pt}

\disc{Coherent velocity wave in the Radcliffe Wave} --- \textit{MNRAS Letters} (2022)\\
Discovered a coherent wave-like velocity pattern across the Galactic-scale Radcliffe structure;
\textbf{first author.}\\
\textbf{G.-X. Li} \& B.-Q. Chen.

% ---------- Education ----------
\vspace{11pt}
\section*{Education}
\textbf{Ph.D. in Astrophysics and Computational Modeling} \hfill \textbf{2014}\\
Max Planck Institute for Radio Astronomy \& University of Bonn, Germany \textit{(magna cum laude)}\\[1pt]
\textbf{M.Sc. in Astrophysics} \hfill \textbf{2010}\\
University of Science and Technology of China (USTC)\\[1pt]
\textbf{B.Sc. in Electronics} \hfill \textbf{2007}\\
Peking University (PKU)

\vspace{6pt}
\section*{Technical Skills}
Methods: Self-supervised learning $\cdot$ representation learning $\cdot$ PyTorch $\cdot$ GPU/HPC
computing $\cdot$ multiscale analysis $\cdot$ hydrodynamics/MHD

\end{document}
